% The layout that resume.tex and cv.tex share. Each one loads it with
% % The layout that resume.tex and cv.tex share. Each one loads it with
% % The layout that resume.tex and cv.tex share. Each one loads it with
% % The layout that resume.tex and cv.tex share. Each one loads it with
% \input{shared}, then adds its own section titles and list style.
% build.py writes build/settings.tex (font, size, margins and spacing) and
% build/body.tex (the content from the JSON file) before each run.

% ---------- Packages ----------
\usepackage{fontspec}
\usepackage{enumitem}
\input{build/settings}
\resumemainfont
\usepackage[margin=\resumemargin]{geometry}
\usepackage[hidelinks]{hyperref}

% ---------- Page layout ----------
\raggedbottom
\pagestyle{empty}
\setlength{\parindent}{0pt}
\setlength{\parskip}{0pt}
\raggedright
% No hyphenated line breaks: a word split across lines can't be found by an
% applicant tracking system's keyword search.
\hyphenpenalty=10000
\exhyphenpenalty=10000

% ---------- Header ----------
% Name and contact lines, centered. The first section title adds the same
% \sectiongap as every other one, so no gap on the page is a one-off.
\newenvironment{resumeheader}{\par\centering\hyphenpenalty=50 \exhyphenpenalty=50 }{\par}
% A separator between header items. If the line wraps there, the bar is
% dropped so no line starts or ends with one. The header turns these breaks
% back on. build.py keeps each header item in a box, so a URL never breaks
% at its own hyphen.
\newcommand{\headersep}{\discretionary{}{}{\hbox{\ \textbar\ }}}

% ---------- Entry macros ----------
\newlength{\datewidth}
% Title on the left (wraps if long), date at the right margin.
\newcommand{\entry}[2]{%
  \settowidth{\datewidth}{#2}%
  \parbox[t]{\dimexpr\linewidth-\datewidth-1em\relax}{\raggedright #1}%
  \hfill #2\par}
% Education: the degree line, then the school with the date at the right
% margin, as eeprep's resume page lays it out.
\newcommand{\eduentry}[3]{%
  {\raggedright #1\par}%
  \entry{#3}{#2}}
% A job's location, on its own line at the right margin under the dates,
% kept with the line above it.
\newcommand{\entrylocation}[1]{\nobreak{\raggedleft #1\par}}
% Space between entries in a section.
\newcommand{\entrygap}{\par\vspace{\entryspace}}

% ---------- Measurements for the fitting script ----------
\newcommand{\measure}[1]{%
  \typeout{RESUME-#1 page=\thepage\space used=\the\pagetotal\space
    goal=\the\pagegoal\space line=\the\baselineskip}}
, then adds its own section titles and list style.
% build.py writes build/settings.tex (font, size, margins and spacing) and
% build/body.tex (the content from the JSON file) before each run.

% ---------- Packages ----------
\usepackage{fontspec}
\usepackage{enumitem}
\input{build/settings}
\resumemainfont
\usepackage[margin=\resumemargin]{geometry}
\usepackage[hidelinks]{hyperref}

% ---------- Page layout ----------
\raggedbottom
\pagestyle{empty}
\setlength{\parindent}{0pt}
\setlength{\parskip}{0pt}
\raggedright
% No hyphenated line breaks: a word split across lines can't be found by an
% applicant tracking system's keyword search.
\hyphenpenalty=10000
\exhyphenpenalty=10000

% ---------- Header ----------
% Name and contact lines, centered. The first section title adds the same
% \sectiongap as every other one, so no gap on the page is a one-off.
\newenvironment{resumeheader}{\par\centering\hyphenpenalty=50 \exhyphenpenalty=50 }{\par}
% A separator between header items. If the line wraps there, the bar is
% dropped so no line starts or ends with one. The header turns these breaks
% back on. build.py keeps each header item in a box, so a URL never breaks
% at its own hyphen.
\newcommand{\headersep}{\discretionary{}{}{\hbox{\ \textbar\ }}}

% ---------- Entry macros ----------
\newlength{\datewidth}
% Title on the left (wraps if long), date at the right margin.
\newcommand{\entry}[2]{%
  \settowidth{\datewidth}{#2}%
  \parbox[t]{\dimexpr\linewidth-\datewidth-1em\relax}{\raggedright #1}%
  \hfill #2\par}
% Education: the degree line, then the school with the date at the right
% margin, as eeprep's resume page lays it out.
\newcommand{\eduentry}[3]{%
  {\raggedright #1\par}%
  \entry{#3}{#2}}
% A job's location, on its own line at the right margin under the dates,
% kept with the line above it.
\newcommand{\entrylocation}[1]{\nobreak{\raggedleft #1\par}}
% Space between entries in a section.
\newcommand{\entrygap}{\par\vspace{\entryspace}}

% ---------- Measurements for the fitting script ----------
\newcommand{\measure}[1]{%
  \typeout{RESUME-#1 page=\thepage\space used=\the\pagetotal\space
    goal=\the\pagegoal\space line=\the\baselineskip}}
, then adds its own section titles and list style.
% build.py writes build/settings.tex (font, size, margins and spacing) and
% build/body.tex (the content from the JSON file) before each run.

% ---------- Packages ----------
\usepackage{fontspec}
\usepackage{enumitem}
\input{build/settings}
\resumemainfont
\usepackage[margin=\resumemargin]{geometry}
\usepackage[hidelinks]{hyperref}

% ---------- Page layout ----------
\raggedbottom
\pagestyle{empty}
\setlength{\parindent}{0pt}
\setlength{\parskip}{0pt}
\raggedright
% No hyphenated line breaks: a word split across lines can't be found by an
% applicant tracking system's keyword search.
\hyphenpenalty=10000
\exhyphenpenalty=10000

% ---------- Header ----------
% Name and contact lines, centered. The first section title adds the same
% \sectiongap as every other one, so no gap on the page is a one-off.
\newenvironment{resumeheader}{\par\centering\hyphenpenalty=50 \exhyphenpenalty=50 }{\par}
% A separator between header items. If the line wraps there, the bar is
% dropped so no line starts or ends with one. The header turns these breaks
% back on. build.py keeps each header item in a box, so a URL never breaks
% at its own hyphen.
\newcommand{\headersep}{\discretionary{}{}{\hbox{\ \textbar\ }}}

% ---------- Entry macros ----------
\newlength{\datewidth}
% Title on the left (wraps if long), date at the right margin.
\newcommand{\entry}[2]{%
  \settowidth{\datewidth}{#2}%
  \parbox[t]{\dimexpr\linewidth-\datewidth-1em\relax}{\raggedright #1}%
  \hfill #2\par}
% Education: the degree line, then the school with the date at the right
% margin, as eeprep's resume page lays it out.
\newcommand{\eduentry}[3]{%
  {\raggedright #1\par}%
  \entry{#3}{#2}}
% A job's location, on its own line at the right margin under the dates,
% kept with the line above it.
\newcommand{\entrylocation}[1]{\nobreak{\raggedleft #1\par}}
% Space between entries in a section.
\newcommand{\entrygap}{\par\vspace{\entryspace}}

% ---------- Measurements for the fitting script ----------
\newcommand{\measure}[1]{%
  \typeout{RESUME-#1 page=\thepage\space used=\the\pagetotal\space
    goal=\the\pagegoal\space line=\the\baselineskip}}
, then adds its own section titles and list style.
% build.py writes build/settings.tex (font, size, margins and spacing) and
% build/body.tex (the content from the JSON file) before each run.

% ---------- Packages ----------
\usepackage{fontspec}
\usepackage{enumitem}
\input{build/settings}
\resumemainfont
\usepackage[margin=\resumemargin]{geometry}
\usepackage[hidelinks]{hyperref}

% ---------- Page layout ----------
\raggedbottom
\pagestyle{empty}
\setlength{\parindent}{0pt}
\setlength{\parskip}{0pt}
\raggedright
% No hyphenated line breaks: a word split across lines can't be found by an
% applicant tracking system's keyword search.
\hyphenpenalty=10000
\exhyphenpenalty=10000

% ---------- Header ----------
% Name and contact lines, centered. The first section title adds the same
% \sectiongap as every other one, so no gap on the page is a one-off.
\newenvironment{resumeheader}{\par\centering\hyphenpenalty=50 \exhyphenpenalty=50 }{\par}
% A separator between header items. If the line wraps there, the bar is
% dropped so no line starts or ends with one. The header turns these breaks
% back on. build.py keeps each header item in a box, so a URL never breaks
% at its own hyphen.
\newcommand{\headersep}{\discretionary{}{}{\hbox{\ \textbar\ }}}

% ---------- Entry macros ----------
\newlength{\datewidth}
% Title on the left (wraps if long), date at the right margin.
\newcommand{\entry}[2]{%
  \settowidth{\datewidth}{#2}%
  \parbox[t]{\dimexpr\linewidth-\datewidth-1em\relax}{\raggedright #1}%
  \hfill #2\par}
% Education: the degree line, then the school with the date at the right
% margin, as eeprep's resume page lays it out.
\newcommand{\eduentry}[3]{%
  {\raggedright #1\par}%
  \entry{#3}{#2}}
% A job's location, on its own line at the right margin under the dates,
% kept with the line above it.
\newcommand{\entrylocation}[1]{\nobreak{\raggedleft #1\par}}
% Space between entries in a section.
\newcommand{\entrygap}{\par\vspace{\entryspace}}

% ---------- Measurements for the fitting script ----------
\newcommand{\measure}[1]{%
  \typeout{RESUME-#1 page=\thepage\space used=\the\pagetotal\space
    goal=\the\pagegoal\space line=\the\baselineskip}}
